\clearpage

\chapter{Expos\'e}
% =====================================================================
% OUTLINE  (max. 2 pages incl. everything except bibliography)
% Budget: ~1 000 words. Follows the structure of the WS exposé:
%   1. Motivation & relevance          (~150 words)
%   2. Problem statement / gap         (~150 words)
%   3. Research question + hypotheses  (~150 words)
%   4. Approach: bridge + model + use cases (~250 words, 1 table)
%   5. Evaluation methodology          (~150 words)
%   6. Work plan / milestones          (~100 words, 1 table)
% Annotated bibliography follows automatically (\nocite{*}).
% =====================================================================

% ---------------------------------------------------------------------
\section*{Motivation}
% - MBSE + SysML v2 (formal KerML semantics, textual notation, standard API/services) = machine-readable models
% - LLM agents + tool use (MCP) have become the dominant integration pattern; MBSE vendors are moving fast
%   (Dassault: SysML v2 REST API on Teamwork Cloud, MagicLab; community MCP servers for Cameo)
% - Industry-standard tool at FH: Magic Systems of Systems Architect (MSoSA) → practical relevance
% Cite: omg-sysml2, omg-sysml2-api, helle2026apollo, mcp-spec
\textit{[TODO ~150 words]}

% ---------------------------------------------------------------------
\section*{Problem Statement}
% - Existing work: NL→SysML v2 generation (prompting, KGs, multi-agent), retrieval (GraphRAG), file/LSP-based
%   validation loops (SEI). All file-based, text-in/text-out, small hand-made models.
% - Nobody has evaluated what an agent can/cannot do against a *live model in an industrial tool*,
%   on a *large* model, across the full task spectrum (query → validate → verify → correct → create).
% - Unknown: capability boundaries; whether tool-side semantics (validation, evaluation) or a thin CRUD
%   bridge is what matters.
% Cite: llm-system-modeling-review, nl2sysml, kg-hallucination, sysforge, graphrag-mbse, sei-native-ai, cae-agents
\textit{[TODO ~150 words]}

% ---------------------------------------------------------------------
\section*{Research Question and Hypotheses}
\textbf{RQ:} \textit{Which SysML v2 modeling tasks can an LLM-based agent perform reliably in an industrial modeling environment (MSoSA) when connected through an MCP bridge, and where do its capabilities break down?}

\begin{itemize}
    \item \textbf{H1 (Task asymmetry):} Read-only tasks (Abfragen, Validierung) are performed with high reliability, whereas write tasks (Korrektur, Erstellen) and Verifizierung degrade with model size and cross-layer dependency depth.
    \item \textbf{H2 (Tool semantics):} Exposing tool-native capabilities (validation suites, expression evaluation) as MCP tools yields larger gains than a thin CRUD wrapper over the SysML v2 API alone. % ablation
    \item \textbf{H3 (Model scale):} Findings obtained on small teaching models (GfSE collection) do not transfer to a large multi-layer model (Apollo 11) without a semantic navigation layer.
\end{itemize}

% ---------------------------------------------------------------------
\section*{Approach}
% Artifact: MCP server bridging an agent to MSoSA
%   Option A: Teamwork Cloud SysML v2 REST API (standard)  |  Option B: MagicDraw OpenAPI Java plugin
%   → decide after checking FH licence (TWC available?). Design both as tool "tiers": thin vs. semantic.
% Model under test: Airbus Apollo 11 SysML v2 model (Helle & Schramm 2026; ~7 kLOC, ~2 000 elements, 5 CoSMA layers)
%   Secondary set: GfSE SysML-v2-Models (small, curated; also "bad examples")
% Use cases (from kickoff):
\begin{table}[h]
\centering
\small
\begin{tabular}{p{2.6cm} p{6.4cm} p{6.4cm}}
\toprule
\textbf{Use case} & \textbf{Example task on Apollo 11} & \textbf{Ground truth / metric} \\
\midrule
Abfragen (query)          & ``Which technical components satisfy requirement X?'' & exact set match vs. model query \\
Validierung (validation)  & ``Is package P well-formed and CoSMA-conformant?'' & tool validation suite results \\
Verifizierung (verification) & ``Is mass budget requirement R satisfied by the current values?'' & expression evaluation in tool \\
Korrektur (correction)    & repair injected faults (dangling ref., wrong type, broken \texttt{satisfy}) & fault-repair rate, no collateral change \\
Erstellen (creation)      & add missing technical component with ports + satisfy links & syntactic validity + expert rubric \\
\bottomrule
\end{tabular}
\caption{Use-case catalogue (to be refined).}
\end{table}
\textit{[TODO ~150 words prose around the table]}

% ---------------------------------------------------------------------
\section*{Evaluation Methodology}
% - Benchmark: ~36 tasks (5 use cases x 3 difficulty levels x 2-3 instances), fully paired across arms
% - Statistical design: per-task pass rate from k=5 repetitions; Friedman + pairwise Wilcoxon (Holm);
%   GLMM (1|task) as robustness check; effect sizes with CIs; pre-registered acceptance criteria; kappa on a 20% subsample
% - Arms (ablation, cf. SEI 3-arm design): (0) no tool, textual file only; (1) thin CRUD bridge; (2) semantic bridge with validation/evaluation tools
% - Metrics: task success, syntactic validity (tool validation), semantic correctness (ground-truth query / expert rubric),
%   collateral damage (diff of model), tool-call count / tokens, failure taxonomy
% - Models: Claude family only (own Anthropic Max subscription) — two capability tiers (e.g. Opus vs. Sonnet)
%   instead of a cross-vendor comparison; state as threat to validity (cf. SEI: single model too)
% - External evaluator outside the agent loop (cf. Pufibara)
\textit{[TODO ~150 words]}

% ---------------------------------------------------------------------
\section*{Work Plan}
% 20 weeks, 01.09.2026 – 19.01.2027; weekly meetings from Oct.
\begin{table}[h]
\centering
\small
\begin{tabular}{l l p{9.5cm}}
\toprule
\textbf{Weeks} & \textbf{Phase} & \textbf{Deliverable} \\
\midrule
1--4   & Setup       & Exposé; import Apollo 11 into MSoSA; MSoSA/TWC API access; literature base \\
5--9   & Bridge      & MCP server (thin tier + semantic tier); use-case task catalogue v1 \\
10--13 & Benchmark   & task set with ground truth; fault injection; evaluation harness \\
14--16 & Experiments & runs across arms and LLMs; failure analysis \\
17--20 & Writing     & thesis text, discussion, submission (19.01.2027) \\
\bottomrule
\end{tabular}
\caption{Milestones.}
\end{table}

\addtocontents{toc}{\vspace{0.8cm}}
